\section{考察}

\subsection{6ヶ月間の構造変化}

6ヶ月間でノード数は37から41へ（+4名），エッジ数は46から67へ（+21本，約46\%増）と
ネットワークが拡大・緻密化した．
これは，既存メンバー間の新規関係形成と，新規メンバーの参加の両方が
寄与していることを示唆する．

\subsection{中心人物の役割}

入次数中心性と媒介中心性の分析から，以下の特徴が明らかになった：

\begin{itemize}
  \item 女性ノード（特に12-2(f)，20-1(f)，23-1(f)）が高い中心性を示した
  \item これらのノードはグループ間を繋ぐ橋渡し役としても機能している
  \item 男性ノードは比較的小規模なクラスタを形成する傾向が見られた
\end{itemize}

\subsection{性別・グループ間パターン}

エッジの性別組み合わせ分析から，以下のパターンが観察された：

\begin{itemize}
  \item 同性間のエッジ（f→f，m→m）が一定数存在
  \item 異性間のエッジも多く，特に女性が中心となる異性間ネットワークが形成されている
  \item グループ間のエッジも存在し，完全に分離されたコミュニティではない
\end{itemize}

\subsection{ネットワークの連結性}

連結成分分析の結果，Beforeネットワークでは6つの連結成分が存在したが，
Afterネットワークでは2つに減少した．
これは，6ヶ月間で孤立していたサブグループがメインのネットワークに統合されたことを示す．
クラスタ係数も0.091から0.177へと約2倍に増加しており，
ネットワーク全体の結束性が高まったことがわかる．

\section{結論}

本レポートでは，Cytoscapeを用いて友人関係ネットワークの6ヶ月間の変化を分析した．
主な知見は以下の通りである：

\begin{enumerate}
  \item ネットワークは6ヶ月間でエッジ数が約46\%増加し，より密な構造へと変化した
  \item 連結成分数が6から2に減少し，孤立グループの統合が進んだ
  \item 特定の女性ノード（12-2(f)，20-1(f)，23-1(f)）がハブおよび橋渡し役として機能
  \item クラスタ係数が約2倍に増加し，ネットワーク全体の結束性が高まった
\end{enumerate}

今後の課題として，時系列でのより詳細な変化追跡，
属性情報を用いた統計的検定，
および他の集団との比較分析が挙げられる．
